\section{Discussion}
\label{sec:discussion}

\subsection{What the Platform Enables}

The comparative questions that motivated this work become configuration rather than engineering.
A study in the style of Jensen et al.~\cite{jensen2025context}, comparing intervention designs on recovery cost, needs no new system: interventions, decision strategies, and analysis providers are selected per session, and every trial leaves a replayable record.
Because decision strategies sit behind a process boundary, machine-learning teams can study their models against live readers without touching, or even reading, the platform's core, as the independent reading-difficulty pipeline in our evaluation demonstrated.
The same seam positions the platform for AI-driven decision strategies, which can be compared against rule-based and human strategies under otherwise identical conditions.

\subsection{Measure the Cost of the Quality You Claim}

The over-repositioning finding generalises beyond this system.
Our restore strategy was plausible, implemented carefully, and wrong: it moved readers three times further than the disturbance it compensated, and the behavioural measures alone would not have told us, because the bundled highlight cue plausibly masked the geometric damage.
The defect surfaced only because the mechanism grades every restore it performs and writes the residual into the record.
We suggest this as a design principle for adaptive interfaces generally: a system that claims a quality attribute, context preservation, smoothness, unobtrusiveness, should instrument the cost of that attribute in production data rather than asserting it.
Two further principles from this work: enforce module boundaries in the build rather than in convention, and treat an outsider's independent implementation, not the authors' own reference implementations, as the test that a contract is truly assumption-free.

\subsection{Limitations}

We state the limits plainly.
The behavioural comparison covers four participants, one text, and manual intervention only; the participants included members of the project team; nothing here is an efficacy claim.
The with-preservation condition bundles positional restore with a highlight cue, so the behavioural gap cannot be attributed to geometry alone, particularly since prior work suggests highlighting itself drives recovery~\cite{jensen2025context}.
The revised restore is verified over a deterministic geometric sweep, not yet with readers.
The automated decision path was validated in a single advisory run rather than a campaign.
The analysis events in the evaluation were produced by the external provider's algorithms, which have not been calibrated against published oculomotor norms.
The webcam facial-state module is integrated but unvalidated, and its difficulty heuristic is hand-tuned.
Finally, the frontend is verified by manual walkthrough rather than an automated suite, and the platform currently renders Markdown content only.
